% Figure 1 (中文版): 评分协议三合一结构图
% 用法：在 paper-zh-tmlr.tex 的 §1 末尾 % Figure 1 (中文版): 评分协议三合一结构图
% 用法：在 paper-zh-tmlr.tex 的 §1 末尾 % Figure 1 (中文版): 评分协议三合一结构图
% 用法：在 paper-zh-tmlr.tex 的 §1 末尾 % Figure 1 (中文版): 评分协议三合一结构图
% 用法：在 paper-zh-tmlr.tex 的 §1 末尾 \input{fig1-scoring-protocol-zh}
% 依赖：\usepackage{tikz}, \usetikzlibrary{positioning,arrows.meta,calc}
% 说明：TikZ 逻辑与英文版 fig1-scoring-protocol.tex 完全一致，仅标签改为中文
\begin{figure}[t]
\centering
\begin{tikzpicture}[
  node distance=0.9cm and 1.1cm,
  box/.style={draw, rounded corners, align=center, minimum height=0.9cm, font=\small},
  lbl/.style={font=\scriptsize},
  arr/.style={->, >=Latex, thick, gray!70}
]
% 左列：分解树
\node[box, draw=blue!60, fill=blue!8] (root) {分解树\\(toy: 2-level)};
\node[box, below=of root] (s) {浅层判断 $S$\\(等速率·占比枚举)};
\node[box, below=of s] (d) {深层判断 $D$\\(真实速率 $t_i=f_i/r_i$)};

% 中列：量
\node[box, right=1.9cm of s, draw=red!60, fill=red!8] (inc) {增量信息\\$H(D|S)=H(D)-I(S;D)$};
\node[box, below=of inc] (score) {可解释分数\\flip $=P(S\neq D)$ · shift};

% 右列：对照
\node[box, right=2.0cm of inc, draw=green!60, fill=green!8] (flat) {平坦速率对照\\all $r_i=1$};
\node[box, below=of flat, draw=green!40, fill=green!5] (zero) {结果：\\$H(D|S)=0$ 精确};
\node[box, right=1.9cm of score, draw=orange!60, fill=orange!8] (perm) {24 参数—深度置换\\(多重集固定)};
\node[box, below=of perm, draw=orange!40, fill=orange!5] (inv) {结果：增量不变\\(去深度化)};

\draw[arr] (root) -- (s);
\draw[arr] (s) -- (d);
\draw[arr] (s.east) -- node[lbl, above] {条件化} (inc.west);
\draw[arr] (d.east) -- (inc.west);
\draw[arr] (inc) -- (score);
\draw[arr] (inc.east) -- node[lbl, above] {施于同一量} (flat.west);
\draw[arr] (flat) -- (zero);
\draw[arr] (score.east) -- (perm.west);
\draw[arr] (perm) -- (inv);
\end{tikzpicture}
\caption{评分协议：分解树 $\rightarrow$ 条件信息增量 $\rightarrow$ 两个负控制。左列给出 toy 结构；中列是闭式可算的评分量；右列两个臂分别排除``深层只是转发浅层''（平坦对照精确归零）与``深层因定义占优''（24 置换不变）。全部判据先于数据生成提交（预注册）。}
\label{fig:scoring-protocol}
\end{figure}

% 依赖：\usepackage{tikz}, \usetikzlibrary{positioning,arrows.meta,calc}
% 说明：TikZ 逻辑与英文版 fig1-scoring-protocol.tex 完全一致，仅标签改为中文
\begin{figure}[t]
\centering
\begin{tikzpicture}[
  node distance=0.9cm and 1.1cm,
  box/.style={draw, rounded corners, align=center, minimum height=0.9cm, font=\small},
  lbl/.style={font=\scriptsize},
  arr/.style={->, >=Latex, thick, gray!70}
]
% 左列：分解树
\node[box, draw=blue!60, fill=blue!8] (root) {分解树\\(toy: 2-level)};
\node[box, below=of root] (s) {浅层判断 $S$\\(等速率·占比枚举)};
\node[box, below=of s] (d) {深层判断 $D$\\(真实速率 $t_i=f_i/r_i$)};

% 中列：量
\node[box, right=1.9cm of s, draw=red!60, fill=red!8] (inc) {增量信息\\$H(D|S)=H(D)-I(S;D)$};
\node[box, below=of inc] (score) {可解释分数\\flip $=P(S\neq D)$ · shift};

% 右列：对照
\node[box, right=2.0cm of inc, draw=green!60, fill=green!8] (flat) {平坦速率对照\\all $r_i=1$};
\node[box, below=of flat, draw=green!40, fill=green!5] (zero) {结果：\\$H(D|S)=0$ 精确};
\node[box, right=1.9cm of score, draw=orange!60, fill=orange!8] (perm) {24 参数—深度置换\\(多重集固定)};
\node[box, below=of perm, draw=orange!40, fill=orange!5] (inv) {结果：增量不变\\(去深度化)};

\draw[arr] (root) -- (s);
\draw[arr] (s) -- (d);
\draw[arr] (s.east) -- node[lbl, above] {条件化} (inc.west);
\draw[arr] (d.east) -- (inc.west);
\draw[arr] (inc) -- (score);
\draw[arr] (inc.east) -- node[lbl, above] {施于同一量} (flat.west);
\draw[arr] (flat) -- (zero);
\draw[arr] (score.east) -- (perm.west);
\draw[arr] (perm) -- (inv);
\end{tikzpicture}
\caption{评分协议：分解树 $\rightarrow$ 条件信息增量 $\rightarrow$ 两个负控制。左列给出 toy 结构；中列是闭式可算的评分量；右列两个臂分别排除``深层只是转发浅层''（平坦对照精确归零）与``深层因定义占优''（24 置换不变）。全部判据先于数据生成提交（预注册）。}
\label{fig:scoring-protocol}
\end{figure}

% 依赖：\usepackage{tikz}, \usetikzlibrary{positioning,arrows.meta,calc}
% 说明：TikZ 逻辑与英文版 fig1-scoring-protocol.tex 完全一致，仅标签改为中文
\begin{figure}[t]
\centering
\begin{tikzpicture}[
  node distance=0.9cm and 1.1cm,
  box/.style={draw, rounded corners, align=center, minimum height=0.9cm, font=\small},
  lbl/.style={font=\scriptsize},
  arr/.style={->, >=Latex, thick, gray!70}
]
% 左列：分解树
\node[box, draw=blue!60, fill=blue!8] (root) {分解树\\(toy: 2-level)};
\node[box, below=of root] (s) {浅层判断 $S$\\(等速率·占比枚举)};
\node[box, below=of s] (d) {深层判断 $D$\\(真实速率 $t_i=f_i/r_i$)};

% 中列：量
\node[box, right=1.9cm of s, draw=red!60, fill=red!8] (inc) {增量信息\\$H(D|S)=H(D)-I(S;D)$};
\node[box, below=of inc] (score) {可解释分数\\flip $=P(S\neq D)$ · shift};

% 右列：对照
\node[box, right=2.0cm of inc, draw=green!60, fill=green!8] (flat) {平坦速率对照\\all $r_i=1$};
\node[box, below=of flat, draw=green!40, fill=green!5] (zero) {结果：\\$H(D|S)=0$ 精确};
\node[box, right=1.9cm of score, draw=orange!60, fill=orange!8] (perm) {24 参数—深度置换\\(多重集固定)};
\node[box, below=of perm, draw=orange!40, fill=orange!5] (inv) {结果：增量不变\\(去深度化)};

\draw[arr] (root) -- (s);
\draw[arr] (s) -- (d);
\draw[arr] (s.east) -- node[lbl, above] {条件化} (inc.west);
\draw[arr] (d.east) -- (inc.west);
\draw[arr] (inc) -- (score);
\draw[arr] (inc.east) -- node[lbl, above] {施于同一量} (flat.west);
\draw[arr] (flat) -- (zero);
\draw[arr] (score.east) -- (perm.west);
\draw[arr] (perm) -- (inv);
\end{tikzpicture}
\caption{评分协议：分解树 $\rightarrow$ 条件信息增量 $\rightarrow$ 两个负控制。左列给出 toy 结构；中列是闭式可算的评分量；右列两个臂分别排除``深层只是转发浅层''（平坦对照精确归零）与``深层因定义占优''（24 置换不变）。全部判据先于数据生成提交（预注册）。}
\label{fig:scoring-protocol}
\end{figure}

% 依赖：\usepackage{tikz}, \usetikzlibrary{positioning,arrows.meta,calc}
% 说明：TikZ 逻辑与英文版 fig1-scoring-protocol.tex 完全一致，仅标签改为中文
\begin{figure}[t]
\centering
\begin{tikzpicture}[
  node distance=0.9cm and 1.1cm,
  box/.style={draw, rounded corners, align=center, minimum height=0.9cm, font=\small},
  lbl/.style={font=\scriptsize},
  arr/.style={->, >=Latex, thick, gray!70}
]
% 左列：分解树
\node[box, draw=blue!60, fill=blue!8] (root) {分解树\\(toy: 2-level)};
\node[box, below=of root] (s) {浅层判断 $S$\\(等速率·占比枚举)};
\node[box, below=of s] (d) {深层判断 $D$\\(真实速率 $t_i=f_i/r_i$)};

% 中列：量
\node[box, right=1.9cm of s, draw=red!60, fill=red!8] (inc) {增量信息\\$H(D|S)=H(D)-I(S;D)$};
\node[box, below=of inc] (score) {可解释分数\\flip $=P(S\neq D)$ · shift};

% 右列：对照
\node[box, right=2.0cm of inc, draw=green!60, fill=green!8] (flat) {平坦速率对照\\all $r_i=1$};
\node[box, below=of flat, draw=green!40, fill=green!5] (zero) {结果：\\$H(D|S)=0$ 精确};
\node[box, right=1.9cm of score, draw=orange!60, fill=orange!8] (perm) {24 参数—深度置换\\(多重集固定)};
\node[box, below=of perm, draw=orange!40, fill=orange!5] (inv) {结果：增量不变\\(去深度化)};

\draw[arr] (root) -- (s);
\draw[arr] (s) -- (d);
\draw[arr] (s.east) -- node[lbl, above] {条件化} (inc.west);
\draw[arr] (d.east) -- (inc.west);
\draw[arr] (inc) -- (score);
\draw[arr] (inc.east) -- node[lbl, above] {施于同一量} (flat.west);
\draw[arr] (flat) -- (zero);
\draw[arr] (score.east) -- (perm.west);
\draw[arr] (perm) -- (inv);
\end{tikzpicture}
\caption{评分协议：分解树 $\rightarrow$ 条件信息增量 $\rightarrow$ 两个负控制。左列给出 toy 结构；中列是闭式可算的评分量；右列两个臂分别排除``深层只是转发浅层''（平坦对照精确归零）与``深层因定义占优''（24 置换不变）。全部判据先于数据生成提交（预注册）。}
\label{fig:scoring-protocol}
\end{figure}
